% frontmatter/english/mapa-del-curso.tex -- el año entero en una tabla.
\section*{The year in this book}

\begin{center}
\renewcommand{\arraystretch}{1.4}
\begin{tabularx}{\linewidth}{@{}c c c c X@{}}
\toprule
\textbf{Part} & \textbf{Days} & \textbf{Season} & \textbf{Words} & \textbf{The story} \\
\midrule
1 & 1--65    & autumn & 30--45 & A big white van stops next door. Who are the new
neighbours? \\
2 & 66--130  & winter & 45--55 & Christmas in two countries, snow, pancakes and a
kite in the wind. \\
3 & 131--195 & spring & 50--65 & A birthday party, an Easter egg hunt, a trip to the
farm -- and where is Pip? \\
4 & 196--260 & summer & 60--75 & A postcard from London, a summer in the village,
and back to school. \\
\bottomrule
\end{tabularx}
\end{center}

\vspace{6mm}
Every part has 13 weeks of 5 days, holidays too. During the year the
texts get longer (the \textit{Words} column), and so do the sentences:
more of them have two parts, joined by \textit{because}, \textit{when},
\textit{if}, \textit{who}, \textit{that}, \textit{before} or
\textit{after}. The story is told in the present in autumn and winter,
and in the past from spring. There are notes, lists, a recipe,
invitations, postcards, letters, a diary and poems.
\newpage
